\documentstyle[% 


righttag% 


,11pt% 


,amssymb% 


,russian%               remove in Eng. 


,izvru%                 Set izven% in Eng. 


% ,izvru-ks             for brief communications 


]{amsart} 


 


%       Math definitions 


 


\newcommand{\comp}{\operatorname{comp}} 


\newcommand{\co}{\operatorname{co}} 


\newcommand{\conv}{\operatorname{conv}} 


\renewcommand{\baselinestretch}{1.05}


 


%%         Use for single entries 


%% \theoremstyle{plain} 


%% \theorembodyfont{\it} 


%% \newtheorem{thmnonum}{\hskip\parindent Теорема} 


%% \renewcommand{\thethmnonum}{} 


 


\theoremstyle{definition} 


\theorembodyfont{\rm} 


\newtheorem{cornonum}{\hskip\parindent Следствие} 


\renewcommand{\thecornonum}{} 


 


 


%\hoffset=-15mm                  For Eng. 


 


\setcounter{page}{13} 


\god{1997} 


\nomer{8~(423)} %                Remove in Eng. 


%\nomer{Vol.\,41, No.\,7}        Set in Eng. 


% \str{75-81}                    Set in Eng. 


\udk{517.95} 


\author{\it А.Н.\,Витюк} 


\title{Существование решений дифференциальных включений\\ 


с частными производными дробных порядков} 


 


\thanks{Работа выполнена при частичной поддержке Международной Соросовской 


Программы поддержки образования в области точных наук на Украине (грант 


APU 061014).} 


\date{} 


% \pagestyle{myheadings}         Set in Eng. 


 


\pagestyle{plain} %              Remove in Eng. 


\begin{document} 


 


% \thispagestyle{plain}          Set in Eng. 


 


\maketitle 


 


% Body text 


 


 


Задача Дарбу для гиперболических дифференциальных уравнений с многозначной 


правой частью (дифференциальных включений) изучалась в работах 


\cite{Daw}--\cite{Vit} и других. 


 


В данной статье рассматриваются вопросы, касающиеся существования и 


некоторых свойств решений задачи Дарбу для дифференциальных включений с 


частными производными дробных порядков. 


 


Основные положения теории дробного интегро-дифференцирования, ее приложения 


в теории интегральных и дифференциальных уравнений и обширная библиография 


работ содержится в монографии \cite{Sam}. 


 


{\bf 1${}^\circ$.} Введем обозначения: $E^n$ --- $n$-мерное евклидово 


пространство с нормой 


$\|\cdot\|$ и нулевым элементом $\theta$; $C(P)$, $AC(P)$, $L(P)$, 


$L_\infty(P)$ --- 


соответственно пространства непрерывных, 


абсолютно непрерывных, суммируемых, измеримых по Лебегу и ограниченных в 


существенном функций, которые определены в области $P$; 


$\comp E^n$ ($\conv E^n$) --- множество 


всех непустых компактных (выпуклых и компактных) подмножеств из $E^n$ с 


метрикой Хаусдорфа $\delta(\cdot,\cdot)$; $\rho(\cdot,\cdot)$ --- расстояние 


между точкой и множеством в $E^n$; $\overline{\co}A$ --- 


замыкание выпуклой оболочки множества $A$; $\Gamma(\cdot)$ --- 


гамма-функция. 


Функция $f:P\to E^n$ называется селектором многозначного отображения 


(м.\,о.) $F:P\to\comp E^n$, если $f(x)\in F(x)$ 


для $x\in P$. 


 


Пусть $G=(0,a]\times(0,b]$, $\overline G=[0,a]\times[0,b]$, 


$0<\alpha,\,\beta\le 1$, $f(x,y)\in L(G)$. 


Выражение 


\begin{equation} 


I_0^\alpha f(x,y)=\frac{1}{\Gamma(\alpha)}\int^x_0 


\frac{f(s,y)}{(x-s)^{1-\alpha}}ds,\quad x>0, 


\end{equation} 


называем \cite{Sam} частным левосторонним интегралом Римана-Лиувилля 


дробного порядка $\alpha$ по переменной $x$, а выражение 


\begin{equation} 


I_0^r f(x,y)=\frac{1}{\Gamma(\alpha)\Gamma(\beta)} 


\int^x_0\int^y_0 


\frac{f(s,t)ds\,dt}{(x-s)^{1-\alpha}(y-t)^{1-\beta}}, 


\quad x>0, \ y>0, 


\end{equation} 


--- смешанным левосторонним дробным интегралом Римана-Лиувилля порядка 


$r=(\alpha,\beta)$. 


Если $f(x,y)\in L(G)$, то \cite{Sam} интегралы 


$I_0^\alpha f(x,y)$, $I_0^r f(x,y)$ определены почти всюду 


(п.\,в.) на $G$ и принадлежат пространству $L(G)$. 


 


Пусть $f_{1-\alpha}(x,y)=I_0^{1-\alpha}f(x,y)$, 


$f_{1-r}(x,y)=I_0^{1-r}f(x,y)$, $1-r=(1-\alpha,1-\beta)$. 


Тогда функции 


\begin{align} 


\begin{split} 


D_0^\alpha f(x,y)=D_x f_{1-\alpha}(x,y)&,\quad 


D_0^r f(x,y)=D^2_{xy}f_{1-r}(x,y)\\ 


\Bigl( D_x=\frac{\partial}{\partial x}&,\quad 


D^2_{xy}=\frac{\partial^2}{\partial x\,\partial y}\Bigr) 


\end{split} 


\end{align} 


называем \cite{Sam} соответственно частной левосторонней производной 


Римана-Лиувилля порядка $\alpha$ по переменной $x$ и смешанной 


левосторонней производной Римана-Лиувилля дробного порядка 


$r=(\alpha,\beta)$. 


Дробное интегрирование обладает полугрупповым свойством \cite{Sam} 


\begin{equation} 


I_0^r I_0^{r_1}f(x,y)=I_0^{r+r_1}f(x,y),\quad 


r_1=(\alpha_1,\beta_1),\quad 0<\alpha_1,\beta_1<1, 


\end{equation} 


где $r+r_1$ --- сумма векторов $r$ и $r_1$. 


 


Заметим, что при $\alpha=\beta=0\;$ $I^0_0f(x,y)=f(x,y)$, а при 


$\alpha=\beta=1\;$ 


$I^1_0f(x,y)=\int\limits^x_0\int\limits^y_0f(s,t)ds\,dt$. 


 


{\bf 2${}^\circ$.} Рассмотрим интегральное уравнение Абеля 


\begin{equation} 


\frac{1}{\Gamma(\alpha)\Gamma(\beta)} 


\int^x_0\int^y_0 


\frac{\varphi(s,t)ds\,dt}{(x-s)^{1-\alpha}(y-t)^{1-\beta}}=f(x,y), 


\quad x>0, \ y>0. 


\end{equation} 


 


\begin{lem} Для того чтобы уравнение {\em (5)} было разрешимо в 


пространстве суммируемых функций, необходимо и достаточно, чтобы 


\begin{align} 


\begin{split} 


&f_{1-r}(x,y)\in AC(G),\\ 


f_{1-r}(x,0)=f_{1-r}&(0,y)=\theta,\quad x\in[0,a], 


\quad y\in[0,b]. 


\end{split} 


\end{align} 


При этом уравнение {\em (5)} имеет единственное решение. 


\end{lem} 


 


Следуя \cite{Sam}, через $I^r_0(L)$ обозначим класс функций $f(x,y)$, 


представимых в виде $f(x,y)=I_0^r\gamma(x,y)$, $\gamma(x,y)\in L(G)$. 


 


\begin{lem} Для того чтобы $f(x,y)\in I_0^r(L)$, необходимо и достаточно, 


чтобы функция $f(x,y)$ удовлетворяла условиям {\em (6)}. 


\end{lem} 


 


Доказательство лемм 1 и 2 аналогично доказательству соответственно теорем 


2.1 и 2.3 из \cite{Sam}. 


 


\begin{lem} Пусть $0<r\le 1$. Тогда для п.\,в. $(x,y)\in G$ 


\medskip 





\noindent {\em 1)} $D_0^rI_0^rf(x,y)=f(x,y)$, 


если $f(x,y)\in L(G);$ \hfill {\em (7)} 


\medskip


 


\noindent {\em 2)} $I_0^rD_0^rf(x,y)=f(x,y)$, если $f(x,y)\in I_0^r(L);$ 


\medskip


 


\noindent {\em 3)} $I^r_0D^r_0f(x,y)=f(x,y)-\mu(x,y)$, 


\hfill {\em (8)}


\medskip


 


\noindent где 


\begin{equation} 


\mu(x,y)=\frac{x^{\alpha-1}}{\Gamma(\alpha)}I_0^\beta\psi'(y)+ 


\frac{y^{\beta-1}}{\Gamma(\beta)}I_0^\alpha\varphi'(x)+ 


\frac{x^{\alpha-1}y^{\beta-1}}{\Gamma(\alpha)\Gamma(\beta)} 


\varphi(0), 


\tag{9} 


\end{equation} 





\begin{gather} 


\text{если }f(x,y)\in L(G),\quad f_{1-r}(x,y)\in AC(G),\notag\\ 


f_{1-r}(x,0)=\varphi(x),\quad x\in[0,a];\quad 


\varphi'(x)\in L_\infty([0,a]),\tag{10}\\ 


f_{1-r}(0,y)=\psi(y),\quad y\in[0,b];\quad 


\psi'(y)\in L_\infty([0,b]),\quad 


\varphi(0)=\psi(0).\notag 


\end{gather} 


\end{lem} 


 


\begin{pf} В силу (3), (4) для п.\,в. $(x,y)\in G$ 


$$ 


D^r_0I^r_0 f(x,y)=D^2_{xy}(I_0^{1-r}I_0^rf(x,y))=f(x,y). 


$$ 


Если $f(x,y)\in I_0^r(L)$, то $f(x,y)=I_0^r\gamma(x,y)$. Тогда 


$$ 


I_0^rD_0^rf(x,y)=I_0^rD^2_{xy}(I_0^{1-r}\gamma(x,y))= 


I_0^r\gamma(x,y)=f(x,y). 


$$ 


Докажем (8). Так как $f_{1-r}(x,y)\in AC(\overline G)$, то п.\,в. на 


$\overline G$ 


существует смешанная левосторонняя 


производная Римана-Лиувилля $D_0^rf(x,y)$, которую обозначим через 


$\lambda(x,y)$. В 


силу условия (10) функцию $f_{1-r}(x,y)$ можно представить в виде 


\cite{Wal} 


\begin{gather*} 


f_{1-r}(x,y)=\omega(x,y)+I_0^1\lambda(x,y),\\ 


\omega(x,y)=\varphi(x)+\psi(y)-\varphi(0). 


\end{gather*} 


Теперь, используя (4), получаем 


\begin{align} 


I_0^{1-r}(f(x,y)-I_0^r\lambda(x,y)) &=\omega(x,y),\notag\\ 


I_0^1(f(x,y)-I_0^r\lambda(x,y)) &=I_0^r\omega(x,y).\tag{11} 


\end{align} 


Так как $f(x,y)-I_0^r\lambda(x,y)\in L(G)$, то из (11) следует, что 


$I_0^r\omega(x,y)\in AC(G)$. 


Следовательно, для п.\,в. $(x,y)\in\overline G$ 


\begin{equation} 


f(x,y)-I_0^r\lambda(x,y)=D^2_{xy}I_0^r\omega(x,y).\tag{12} 


\end{equation} 


Докажем, что $D^2_{xy}I_0^r\omega(x,y)$ совпадает с $\mu(x,y)$. 


Пусть $v(x,y)=I_0^r\varphi(x)$. Так 


как $\varphi(x)=\varphi(0)+$\linebreak %%%%%%%%% REMOVE!!


$+\int\limits_0^x\varphi'(t)dt$, то, используя 


теорему Фубини, получаем 


\begin{equation} 


v(x,y)=\frac{x^\alpha y^\beta\varphi(0)} 


{\Gamma(1+\alpha)\Gamma(1+\beta)}+ 


\frac{y^\beta}{\Gamma(1+\alpha)\Gamma(1+\beta)}\int_0^x 


(x-s)^\alpha\varphi'(s)ds.\tag{12${}_1$} 


\end{equation} 


Тогда 


\begin{equation} 


D^2_{xy}v(x,y)= 


\frac{x^{\alpha-1}y^{\beta-1}\varphi(0)}{\Gamma(\alpha)\Gamma(\beta)} 


+\frac{y^{\beta-1}}{\Gamma(\beta)}I_0^\alpha\varphi'(x),\tag{13} 


\end{equation} 


причем $I_0^\alpha\varphi'(x)\in C([0,a])$. Аналогично доказываем, что 


\begin{equation} 


D^2_{xy}I_0^r\psi(y)= 


\frac{x^{\alpha-1}y^{\beta-1}\varphi(0)}{\Gamma(\alpha)\Gamma(\beta)} 


+\frac{x^{\alpha-1}}{\Gamma(\alpha)}I_0^\beta\psi'(y).\tag{14} 


\end{equation} 


Из соотношений (12)--(14) следует (8), причем $\mu(x,y)\in C(G)$. 


\end{pf} 


 


{\bf 3${}^\circ$.} Рассмотрим дифференциальное включение (д.\,в.) 


\begin{equation} 


D_0^ru(x,y)\in F(x,y,u(x,y))\tag{15} 


\end{equation} 


и краевые условия 


\begin{align} 


u_{1-r}(x,0) &=\varphi(x),\quad x\in[0,a],\notag\\ 


u_{1-r}(0,y) &=\psi(y),\quad y\in[0,b],\quad \varphi(0)=\psi(0).\tag{16} 


\end{align} 


 


Предположим, что м.\,о. $F:\overline G\times E^n\to\comp E^n$ удовлетворяет 


условиям: 


\begin{itemize} 


\item[(A${}_1$)] $F(\cdot,\cdot,u):\overline G\to\comp E^n$ измеримо для 


каждого $u\in E^n$; 


\item[(A${}_2$)] $\delta(F(x,y,u),F(x,y,v))\le K\|u-v\|$ для п.\,в. 


$(x,y)\in\overline G$ и 


любых $u,v\in E^n$; 


\item[(A${}_3$)] $|F|=\delta(F(x,y,u),\{\theta\})\le M$ для любых 


$(x,y,u)\in\overline G\times E^n$. 


\end{itemize} 


Предположим также, что функции $\varphi:[0,a]\to E^n$, 


$\psi:[0,b]\to E^n$ абсолютно 


непрерывные и $\varphi'\in L_\infty([0,a])$, 


$\psi'\in L_\infty([0,b])$. 


 


\begin{defn} Функцию $u:G\to E^n$ называем решением задачи (15)--(16), 


если 


\begin{itemize} 


\item[(а)] $u(x,y)\in C(G)$; 


\item[(б)] $u_{1-r}(x,y)\in AC(\overline G)$; 


\item[(в)] $u(x,y)$ удовлетворяет краевым условиям (16) и п.\,в. на $G$ 


д.\,в. (15). 


\end{itemize} 


\end{defn} 


 


Рассмотрим дифференциальное уравнение 


\begin{equation} 


D_0^ru(x,y)=\gamma(x,y).\tag{17} 


\end{equation} 


 


\begin{lem} Если $\gamma(x,y)\in L(G)$ и $\|\gamma(x,y)\|\le M$, то задача 


{\em (17), (16)} эквивалентна уравнению 


\begin{equation} 


u(x,y)=I_0^r\gamma(x,y)+\mu(x,y). \tag{18}


\end{equation} 


\end{lem} 


 


\begin{pf} Если решение задачи (17), (16) существует, то в силу (8), (9) 


его можно записать в виде (18). 


 


Докажем далее, что функция $u(x,y)$ является решением задачи (17), (16). 


В лемме 3 отмечалось, что $\mu(x,y)\in C(G)$. Так как 


$\|\gamma(x,y)\|\le M$, то непосредственно 


проверяем, что 


$I_0^r\gamma(x,y)\in C(G)$. Следовательно, $u(x,y)$ удовлетворяет условию 


(а) определения 1. Остается 


проверить, что $u(x,y)$ удовлетворяет условиям (б), (в) этого определения. 


Пусть $q(x,y)=I_0^r\omega(x,y)$. Аналогично выводу соотношения (12${}_1$) 


получим преобразованное выражение 


$$ 


q(x,y)= 


\frac{1}{\Gamma(1+\alpha)\Gamma(1+\beta)}\biggl[ 


x^\alpha y^\beta\varphi(0)+y^\beta 


\int^x_0(x-s)^\alpha\varphi'(s)ds+x^\alpha 


\int^y_0(y-t)^\beta\psi'(t)dt\biggr], 


$$ 


из которого следует, что для $x\in[0,a]$, $y\in[0,b]$ 


\begin{equation} 


q(x,0)=q(0,y)=\theta.\tag{19} 


\end{equation} 


Согласно лемме 3 $\mu(x,y)=D^2_{xy}q(x,y)$. Выразив $\mu(x,y)$ из (18) и 


вычислив $I_0^1\mu(x,y)$ с учетом (19), получим 


$$ 


I_0^1(u(x,y)-I_0^r\gamma(x,y))=I_0^r\omega(x,y). 


$$ 


Отсюда, используя (4), имеем 


$$ 


I_0^r(I_0^{1-r}u(x,y)-I_0^1\gamma(x,y)-\omega(x,y))=\theta, 


$$ 


поэтому в силу леммы 1\; 


$I_0^{1-r}u(x,y)-I_0^1\gamma(x,y)-\omega(x,y)=\theta$, т.\,е. 


$u_{1-r}(x,y)=I_0^1\gamma(x,y)+\omega(x,y)$. 


Следовательно, $u_{1-r}(x,y)\in AC(\overline G)$ и удовлетворяет условиям 


(16), а также 


$D_0^ru(x,y)=\gamma(x,y)$ для п.\,в. $(x,y)\in G$. 


\end{pf} 


 


\begin{thm} Пусть м.\,о. $F:\overline G\times E^n\to\comp E^n$ 


удовлетворяет условиям {\em (A${}_1$)--(A${}_3$)}, а функция 


$\tau(x,y):G\to E^n$ 


удовлетворяет условиям {\em (а), (б)} определения 


{\em 1}, условиям {\em (16)} и для 


п.в. $(x,y)\in G$ 


\begin{equation} 


\rho(D_0^r\tau(x,y),F(x,y,\tau(x,y)))\le\lambda<+\infty.\tag{20} 


\end{equation} 


 


Тогда существует такое решение $u(x,y)$ задачи 


{\em (15), (16)}, что для $(x,y)\in G$ 


\begin{equation} 


\|u(x,y)-\tau(x,y)\|\le\frac{\lambda}{K}(\widetilde E_r 


(Kx^\alpha y^\beta)-1),\tag{21} 


\end{equation} 


где 


$$ 


\widetilde E_r(z)=\sum_{k=0}^\infty 


\frac{z^k}{\Gamma(\alpha k+1)\Gamma(\beta k+1)}. 


$$ 


\end{thm} 


 


\begin{pf} Построим последовательности функций 


$u^{(k)}(x,y)$, $v_k(x,y)$,\linebreak %%%%%%%% REMOVE


$k\ge 0$, 


$u^{(0)}(x,y)=\tau(x,y)$, $v^{(0)}(x,y)=D_0^r\tau(x,y)$, 


$u^{(k)}(x,y)=I_0^rv_k(x,y)+\mu(x,y)$, $k\ge 1$, 


где $v_k(x,y)$ --- такой измеримый селектор измеримого м.\,о. 


$F(x,y,u^{(k-1)}(x,y))$, что для 


п.\,в. $(x,y)\in G$ 


\begin{equation} 


\|v_k(x,y)-v_{k-1}(x,y)\|= 


\rho(v_{k-1}(x,y), F(x,y,u^{(k-1)}(x,y))).\tag{22} 


\end{equation} 


По индукции получаем 


\begin{align} 


\|v_{m+1}(x,y)-v_m(x,y)\| &\le 


\frac{\lambda(Kx^\alpha y^\beta)^m} 


{\Gamma(m\alpha+1)\Gamma(m\beta+1)} 


\text{ для п.\,в. } (x,y)\in G,\tag{23}\\ 


% 


\|u_{m+1}(x,y)-u_m(x,y)\| &\le \frac{\lambda}{K} 


\frac{(Kx^\alpha y^\beta)^m} 


{\Gamma(m\alpha+1)\Gamma(m\beta+1)},\quad 


(x,y)\in G.\tag{24}


\end{align} 


 


При помощи оценок (23), (24) доказываем, что последовательность 


$v_k(x,y)$, $k\ge 0$, п.\,в. 


на $G$ сходится к функции $v(x,y)\in L(G)$, а последовательность 


$u_k(x,y)$, $k\ge 0$, 


сходится равномерно 


на $G$ к функции $u(x,y)\in C(G)$. Согласно теореме Лебега 


$u(x,y)=I_0^rv(x,y)+\mu(x,y)$, а согласно 


лемме 4 $u(x,y)$ 


удовлетворяет условию (16), 


$u_{1-r}(x,y)\in AC(\overline G)$, 


$D_0^ru(x,y)=v(x,y)$ для п.\,в. $(x,y)\in G$. Из (22) при $k\to\infty$ 


следует, что $u(x,y)$ удовлетворяет п.\,в. на $G$ д.\,в. (15). 


 


Воспользовавшись оценкой (24), получаем, что для $(x,y)\in G$ имеет место 


оценка 


$$ 


\|u^{(m+1)}(x,y)-\tau(x,y)\|\le\sum^m_{k=0} 


\|u^{(k+1)}(x,y)-u^{(k)}(x,y)\|\le 


\frac{\lambda}{K}(\widetilde E_r(Kx^\alpha y^\beta)-1). 


$$ 


Отсюда при $m\to\infty$ следует (21). 


\end{pf} 


 


\begin{cornonum} Пусть 


$\tau(x,y)=D^2_{xy}I_0^r\omega(x,y)$. Тогда, как и в лемме 4, докажем, что 


$\tau_{1-r}(x,y)=\omega(x,y)$. Следовательно, 


$\tau(x,y)$ удовлетворяет условиям (а), (б) определения 1 и краевым 


условиям (16), а 


$D_0^r\tau(x,y)=D^2_{xy}\omega(x,y)=\theta$ для п.\,в. 


$(x,y)\in G$. Поэтому с учетом 


условия (A${}_3$) $\rho(D_0^r\tau(x,y),F(x,y,\tau(x,y)))\le M$ 


и согласно теореме 1 


множество решений задачи (15), (16) непусто. 


\end{cornonum} 


 


{\bf 4${}^\circ$.} Рассмотрим далее д.\,в. (15) и краевые условия 


\begin{equation} 


u_{1-r}(x,0)=u_{1-r}(0,y)=\theta,\quad x\in[0,a],\quad 


y\in [0,b].\tag{25} 


\end{equation} 


Функция $u:\overline G\to E^n$ называется решением задачи (15), (25), если 


$u(x,y)$ удовлетворяет условиям определения 1 в области 


$\overline G$. Если $u(x,y)$ есть 


некоторое 


решение задачи (15), (25), то 


$u(x,y)=I_0^rv(x,y)$, где $v(x,y)$ --- измеримый селектор м.\,о. 


$F(x,y,u(x,y))$. 


Так как $\|v(x,y)\|\le M$, то легко доказать, что 


$\lim I_0^rv(x,y)=\theta$ при $x\to 0$ для каждого $y\in(0,b]$ и 


при $y\to 0$ для 


каждого $x\in(0,a]$. Полагаем $I_0^rv(x,y)=\theta$ при $x=0$, 


$y\in[0,b]$ и $y=0$, $x\in[0,a]$. 


 


Имеет место утверждение, аналогичное теореме 1, причем оценка (21) верна 


для $(x,y)\in\overline G$. Если в следствии положить $\tau(x,y)=\theta$, 


то получим, что множество решений 


задачи (15), (25), которое обозначим через $H(F;r)$, непусто и согласно 


лемме 2 $H(F;r)\subset I_0^r(L)$. 


 


\begin{thm} Пусть м.\,о. $F:\overline G\times E^n\to\conv E^n$ 


удовлетворяет условию {\em (А${}_3$)} и условиям: 


\begin{itemize} 


\item[(B${}_1$)] $F(x,y,\cdot)$ полунепрерывно сверху для п.\,в. 


$(x,y)\in\overline G$; 


\item[(B${}_2$)] для любого $u\in E^n$ существует измеримый селектор 


м.\,о. $F(\cdot,\cdot,u)$. 


\end{itemize} 


Тогда существует решение задачи {\em (15), (25).} 


\end{thm} 


 


\begin{pf} Область $G$ покроем сеткой с узлами 


$(x_i,y_j)$, $x_i=ih$, $y_j=jl$, $mh=a$, $ml=b$, и пусть 


$G_{ij}=[x_i,x_{i+1}]\times[y_j,y_{j+1}]$. 


Построим функцию $u^{(m)}:\overline G\to E^n$, удовлетворяющую краевым 


условиям (25). Пусть 


$\gamma_{00}:G_{00}\to E^n$ --- некоторый измеримый селектор м.\,о. 


$F(x,y,\theta)$. 


Полагаем $u^{(m)}(x,y)=I_0^r\gamma_{00}(x,y)$ в области $G_{00}$. 


Предположим, что функция $U^{(m)}(x,y)$ построена в области 


$[0,x_{i+1}]\times[0,y_{j+1}]\setminus(x_i,x_{i+1}]\times(y_j,y_{j+1}]$ ,


и определим ее в области 


$G_{ij}$. Если $\gamma_{ij}:G_{ij}\to E^n$ --- измеримый селектор м.\,о. 


$F(x,y,u^{(m)}(x_i,y_j))$, 


а $v_m(x,y)=\gamma_{k\nu}(x,y)$ для $(x,y)\in G_{k\nu}$, 


$k=\overline{0,i}$, $\nu=\overline{0,j}$, 


то для $(x,y)\in G_{ij}$ полагаем 


$u^{(m)}(x,y)=I_0^rv_m(x,y)$. Так как $v_m\in L(G)$, то 


$u^{(m)}_{1-r}(x,y)=I_0^1v_m(x,y)\in AC(\overline G)$ и множество 


функций 


$u^{(m)}_{1-r}(x,y)$, $m\ge 1$, является равномерно ограниченным и 


равностепенно непрерывным. 


 


Кроме того, для любых $(x_1,y_1),(x_2,y_2)\in G$, $x_1\le x_2$, 


$y_1\le y_2$, $m\ge 1$ 


\begin{gather} 


\|u^{(m)}(x_2,y_2)-u^{(m)}(x_1,y_1)\|\le 


\frac{M} 


{\Gamma(1+\alpha)\Gamma(1+\beta)}a^\alpha(y_2^\beta-y_1^\beta)+\notag \\


% 


+b^\beta(x_2^\alpha-x_1^\alpha)+2(b^\beta(x_2-x_1)^\alpha+ 


a^\alpha(y_2-y_1)^\beta)+ 


(x_2-x_1)^\alpha(y_2-y_1)^\beta],\tag{26} 


\end{gather} 


а также 


\begin{equation} 


\|u^{(m)}(x,y)\|\le 


\frac{Ma^\alpha b^\beta}{\Gamma(1+\alpha)\Gamma(1+\beta)}, 


\quad m\ge 1. \tag{27}


\end{equation} 


Следовательно, и множество функций $u^{(m)}(x,y)$, $m\ge 1$, является 


равномерно ограниченным 


и равностепенно непрерывным. В соответствии с теоремой Арцела существуют 


подпоследовательности (обозначим их также 


$u^{(m)}(x,y)$ и $u^{(m)}_{1-r}(x,y)$, $m\ge 1$), которые равномерно 


на $\overline G$ сходятся соответственно к $u(x,y)$ и $z(x,y)$. Определим 


функцию $g^{(m)}(x,y):\overline G\to E^n$, полагая 


$g^{(m)}(x,y)=u^{(m)}(x_i,y_j)$, если 


$(x,y)\in[x_i,x_{i+1})\times[y_j,y_{j+1})$, 


$i,j=\overline{0,m-2}$, 


$g^{(m)}(x,y)=u^{(m)}(x_{m-1},y_j)$, 


если $(x,y)\in G_{m-1,j}$, $j=\overline{0,m-1}$, и 


$g^{(m)}(x,y)=u^{(m)}(x_i,y_{m-1})$, 


если $(x,y)\in G_{i,m-1}$, $i=\overline{0,m-2}$. Очевидно, 


$g^{(m)}(x,y)$ также 


равномерно на $\overline G$ сходится к $u(x,y)$ при $m\to\infty$ и 


\begin{equation} 


v_m(x,y)\in F(x,y,g^{(m)}(x,y)).\tag{28} 


\end{equation} 


Воспользовавшись (28) и леммой 1.2 (\cite{Sos}, с.\,19), 


получаем для п.\,в. $(x,y)\in\overline G$ 


\begin{align*} 


D^2_{xy}z(x,y)&\in\bigcap\limits_{i=1}^\infty\overline{\co} 


\bigcup\limits_{k=i}^\infty D^2_{xy}u^{(k)}_{1-r}(x,y)\subset\\ 


% 


&\subset\bigcap\limits_{i=1}^\infty\overline{\co} 


\bigcup\limits_{k=i}^\infty F(x,y,g^{(k)}(x,y))\subset 


F(x,y,u(x,y)). 


\end{align*} 


Поскольку $u^{(m)}_{1-r}(x,y)=I_0^{1-r}u^{(m)}(x,y)$, то в соответствии с 


теоремой Лебега $z(x,y)=I_0^{1-r}u(x,y)=u_{1-r}(x,y)$. Отсюда и из 


(3) следует, что п.\,в. на $\overline G\;$ 


$D_0^ru(x,y)=D^2_{xy}z(x,y)\in F(x,y,u(x,y))$. Этим завершается 


доказательство теоремы 2. 


\end{pf} 


 


\begin{thm} Пусть м.\,о. $F:\overline G\times E^n\to\conv E^n$ 


удовлетворяет условиям {\em (B${}_1$), (B${}_2$)} и существуют такие 


постоянные $R$ и $T$, что 


$|F(x,y,u)|\le R+T\|u\|$ для 


$(x,y,u)\in\overline G\times E^n$. Тогда существует решение задачи 


{\em (15), (25).} 


\end{thm} 


 


Для доказательства теоремы 3 можно воспользоваться техникой доказательства 


теоремы 1 работы \cite{Bul}, адаптированной применительно к задаче (15), 


(25). 


При этом вместо леммы 4 этой работы следует воспользоваться следующим 


аналогом неравенства Вендроффа для интеграла Римана-Лиувилля дробного 


порядка $r$. 


 


\begin{lem} Пусть $v:\overline G\to[0,+\infty)$, 


$v\in C(G)$, и для $(x,y)\in\overline G$ 


$$ 


v(x,y)\le A+\frac{B}{\Gamma(\alpha)\Gamma(\beta)} 


\int^x_0\int^y_0 (x-s)^{\alpha-1}(y-t)^{\beta-1} 


v(s,t)ds\,dt, 


$$ 


где $A$ и $B$ --- неотрицательные постоянные. Тогда для $(x,y)\in\overline G$ 


$$ 


v(x,y)\le A\cdot\widetilde E_r(Bx^\alpha y^\beta). 


$$ 


\end{lem} 


 


\begin{thm} Если м.\,о. $F:\overline G\times E^n\to\conv E^n$ 


удовлетворяет условиям теоремы 


{\em 2}, то множество $H(F;r)$ является компактным подмножеством 


пространства $C(\overline G)$. 


\end{thm} 


 


\begin{pf} Для любого решения $u(x,y)\in H(F;r)$ справедливы оценки (26), 


(27). Следовательно, 


множество $H(F;r)$ является равномерно ограниченным и равностепенно 


непрерывным 


и остается еще доказать, что $H(F;r)$ --- замкнутое множество. Рассмотрим 


последовательность 


$u^{(k)}(x,y)\in H(F;r)$, $k\ge 1$, сходящуюся к 


$u(x,y)\in C(G)$. Так как 


$u^{(k)}(x,y)=I_0^rv_k(x,y)$, где $v_k(x,y)$ --- 


некоторый 


измеримый селектор м.\,о. 


$F(x,y,u^{(k)}(x,y))$, то 


$u^{(k)}_{1-r}(x,y)=I_0^1v_k(x,y)\in AC(\overline G)$. Далее, рассуждая, 


как в теореме 2 при 


обосновании сходимости последовательности 


$u^{(m)}(x,y)$, $m\ge 1$, к решению задачи (15), 


(25), докажем, что $u(x,y)\in H(F;r)$. 


\end{pf}


 


\begin{thebibliography}{9} 


 


\bibitem{Daw} 


Dawidowski M., Kubiaczyk I. {\em On bounded solutions of hyperbolic 


differential inclusion in Banach spaces} // Demonstratio Math. -- 1992. -- 


V.\,25. -- \No\,1--2. -- P.\,153--159. 


 


\bibitem{Sos} 


Sosulski W. {\em W\l{}asno\'sci rozwi\c{a}nza\'n uog\'olnionych 


r\'owna\'n r\'o\.{z}niczkowych o pochodnych cz\c{a}stkowych typu 


hiperbolicznego}. -- Zielona G\'ora, 1982. -- 68 s. 


 


\bibitem{Vas} 


Vasile Staicu. {\em On a non-convex hyperbolic differential inclusion} // 


Proc. Edinburg Math. Soc. -- 1992. -- V.\,35. -- P.\,375--382. 


 


\bibitem{Vit} 


Витюк А.Н. {\em О существовании решений одного класса многозначных 


дифференциальных 


уравнений с частными производными} // Укр. матем. журн. -- 1990. -- 


Т.42. -- \No\,11. -- С.1454--1460. 


 


\bibitem{Sam} 


Самко С.Г., Килбас А.А., Маричев О.И. {\em Интегралы и производные дробного 


порядка и некоторые их приложения}. -- Минск: Наука, 1987. -- 687 с. 


 


\bibitem{Wal} 


Walczak S. {\em Absolutely continuous functions of several variables and 


their application to differential equation} // Bull. Acad. polon. sci. Ser. 


sci. math., astron. et phys. -- 1987. -- V.\,35. -- \No\,11--12. -- P.\,733 


--734. 


 


\bibitem{Bul} 


Булгаков А.И. {\em К вопросу о свойствах множества решений дифференциальных 


включений} // Дифференц. уравнения. -- 1976. - Т.\,12. -- \No\,6. -- 


С.971--977. 


 


\bibitem{Bec} 


Беккенбах Э., Беллман Р. {\em Неравенства}. -- М.: Мир, 1965. -- 276 с. 


 


\end{thebibliography} 


 


\vskip 2cm 


 


\post{\it Одесский университет {\em (}Украина{\em )}}{\it Поступила} 


\post{}{06.02.1995} 


 


\end{document} 


